% © 2026 Ing. David Jaroš — CC BY-NC-ND 4.0
%
% This work is licensed under a Creative Commons Attribution-NonCommercial-NoDerivatives
% 4.0 International License (CC BY-NC-ND 4.0).
%
% License History: Earlier drafts (up to v0.3) were released under CC BY 4.0.
% From v0.4 onward, all material is released under CC BY-NC-ND 4.0 to protect
% the integrity of the theoretical work during ongoing academic development.
%
% See LICENSE.md for full license text.

% UBT-AI-PROVENANCE-BEGIN
% schema: ubt-ai-provenance/v1
% tier: C_working
% ai_assistance: disclosed
% human_review: risk-based
% editorial_responsibility: Ing. David Jaroš
% policy: ../../AI_PROVENANCE.md
% notice: Working material; exhaustive human review is not claimed.
% UBT-AI-PROVENANCE-END

\ifdefined\INCLUDEMODE
\else
  \documentclass[12pt,a4paper]{article}
  \usepackage[a4paper, margin=2.5cm]{geometry}
  \usepackage{amsmath, amssymb, amsthm}
  \usepackage{mathtools}
  \usepackage{hyperref}
  \usepackage{xcolor}
  \usepackage{booktabs}
  \usepackage{tcolorbox}
  \tcbuselibrary{skins,breakable}
  \usepackage{mdframed}

  \newtheorem{theorem}{Theorem}[section]
  \newtheorem{lemma}[theorem]{Lemma}
  \newtheorem{proposition}[theorem]{Proposition}
  \theoremstyle{remark}
  \newtheorem{remark}[theorem]{Remark}

  \newmdenv[backgroundcolor=yellow!15, linecolor=orange!70, linewidth=1.5pt]{gapbox}
  \newmdenv[backgroundcolor=green!12, linecolor=green!60!black, linewidth=1.5pt]{resultbox}

  \newtcolorbox{statusbox}[1][]{colback=gray!8!white,colframe=black!60,
    arc=3pt,boxrule=1pt,fonttitle=\bfseries,title=Status Summary,#1}

  \newcommand{\BB}{\mathbb{B}}
  \newcommand{\CC}{\mathbb{C}}
  \newcommand{\RR}{\mathbb{R}}
  \newcommand{\HH}{\mathbb{H}}
  \newcommand{\ZZ}{\mathbb{Z}}
  \newcommand{\Tr}{\mathrm{Tr}}

  \title{\textbf{Central Charge of the UBT Alpha-Sector CFT}\\[4pt]
         \large Gap G-c3: Derivation Attempt and Status}
  \author{Ing.\ David Jaro\v{s}}
  \date{2026-05-10}

  \begin{document}
  \maketitle
\fi

%──────────────────────────────────────────────────────────────────────────────
\section{Objective}
\label{sec:goal}
%──────────────────────────────────────────────────────────────────────────────

The formula
\begin{equation}
  B = N_{\mathrm{eff}}^{3/2}\,(2\eta(i))^{c/12}
  \label{eq:B_formula}
\end{equation}
requires the Virasoro central charge $c$ of the 2D CFT associated with the
UBT alpha sector.  For the exponent $c/12 = 1/4$, one needs $c = 3$.

This document attempts to derive $c$ from the second functional derivative
$\delta^2 S[\Theta]/\delta\Theta^2$ evaluated at the winding saddle, counts the
physical degrees of freedom of the fluctuation operator, and gives the
resulting verdict on Gap~G-c3.

\textbf{Hard rules}: No use of $\alpha_{\exp}$, $137$, or $B_{\mathrm{req}}$
as inputs.  All input from $S[\Theta]$ and the algebra $\BB = \CC\otimes\HH$.

%──────────────────────────────────────────────────────────────────────────────
\section{Step 1 — Winding Saddle and Fluctuation Operator}
\label{sec:saddle}
%──────────────────────────────────────────────────────────────────────────────

\subsection{Winding background}

The relevant saddle configuration for the alpha-sector winding potential is
(from \texttt{canonical/alpha/alpha\_best\_route.tex} and
\texttt{research\_tracks/T3\_ALPHA/cw\_determinant\_full\_derivation.tex}):
\begin{equation}
  \bar\Theta_n(x,\psi) \;=\; \Theta_0(x)\,e^{in\psi},
  \qquad n \in \ZZ,
  \label{eq:winding_saddle}
\end{equation}
where $\psi \in [0,2\pi)$ is the compact EM direction and $\Theta_0(x)$ is a
slowly varying 4D field.

\subsection{Second functional derivative}

Expanding $\Theta = \bar\Theta_n + \delta\Theta$ around the saddle:
\begin{align}
  S[\bar\Theta_n + \delta\Theta]
  &= S[\bar\Theta_n]
    + \frac{\delta S}{\delta\Theta}\bigg|_{\bar\Theta_n}\!\!\!\delta\Theta
    + \frac{1}{2}
      \int\!\int \delta\Theta^\dagger(y)\,
      \frac{\delta^2 S}{\delta\Theta^\dagger(y)\,\delta\Theta(x)}\bigg|_{\bar\Theta_n}\!\!\!
      \delta\Theta(x)\,dx\,dy + \ldots
  \label{eq:expansion}
\end{align}
From the minimal kinetic action
$S_{\mathrm{kin}}[\Theta] = \int d^4x\,d\psi\;\mathrm{Sc}[(D_\mu\Theta)^\dagger(D^\mu\Theta)]$,
the quadratic fluctuation operator is
\begin{equation}
  \hat{K}_n \;:=\; -D^\mu D_\mu - \partial_\psi^2 + \text{(mass terms from curvature)},
  \label{eq:K_op}
\end{equation}
where $D_\mu = \partial_\mu + ieA_\mu$ is the 4D covariant derivative and the
background field $\bar\Theta_n$ contributes through the winding mass
$m_n^2 = n^2/R_\psi^2$ from the compact $\psi$-direction:
\begin{equation}
  -\partial_\psi^2\,e^{in\psi}\,\delta\Theta
  \;\longrightarrow\; m_n^2\,\delta\Theta.
  \label{eq:winding_mass}
\end{equation}
Therefore on the winding saddle, the fluctuation operator acting on each
real component of $\delta\Theta$ is:
\begin{equation}
  \hat{K}_n \;=\; -\partial^2_\mu + m_n^2,
  \qquad m_n^2 = \frac{n^2}{R_\psi^2}.
  \label{eq:K_final}
\end{equation}
This is the Klein-Gordon operator for a massive scalar of mass $m_n$.

%──────────────────────────────────────────────────────────────────────────────
\section{Step 2 — Identification as a 2D CFT}
\label{sec:CFT}
%──────────────────────────────────────────────────────────────────────────────

\subsection{Reduction to the compact directions}

The one-loop effective action from the fluctuation operator is:
\begin{equation}
  W_{\mathrm{1loop}} \;=\; \frac{1}{2}\,\ln\det\hat{K}_n
  \;=\; \frac{1}{2}\sum_{\text{modes}} \ln E_{\text{mode}}.
  \label{eq:W_1loop}
\end{equation}
At the compactification scale $R_\psi \sim 1/\mu$, the relevant effective
field theory is the restriction of $\delta\Theta$ to the compact $S^1_\psi$
direction.  The resulting 2D theory lives on $S^1_t \times S^1_\psi$
(thermal circle times compact EM circle) and has a 2D kinetic term of the form:
\begin{equation}
  \mathcal{L}_{\mathrm{2D}} \;=\; \frac{1}{2}(\partial_t\varphi)^2
  + \frac{1}{2}(\partial_\psi\varphi)^2
  \label{eq:2D_kin}
\end{equation}
for each real component $\varphi$ of $\delta\Theta$, in the massless limit
$m_n \to 0$ (or at scales $\mu \gg m_n$).

The 2D CFT is therefore a free compact scalar on $T^2 = S^1_t \times S^1_\psi$
with central charge $c = 1$ per real scalar degree of freedom.

\subsection{Standard CFT central-charge assignment}

For $N_{\mathrm{dof}}$ free real scalars on a 2D surface:
\begin{equation}
  c \;=\; N_{\mathrm{dof}}.
  \label{eq:c_formula}
\end{equation}
This is the standard Virasoro central charge from the OPE of the energy-momentum
tensor $T(z)T(0) \sim c/(2z^4) + \ldots$ (Di Francesco et al.,
\emph{Conformal Field Theory}, Springer 1997, §5.2).

%──────────────────────────────────────────────────────────────────────────────
\section{Step 3 — Counting Physical Degrees of Freedom}
\label{sec:dof}
%──────────────────────────────────────────────────────────────────────────────

\subsection{Total real degrees of freedom of $\Theta$}

The biquaternion field $\Theta \in \BB = \CC\otimes\HH$ has:
\begin{equation}
  \dim_\RR(\BB) = \dim_\RR(\CC)\times\dim_\RR(\HH) = 2\times 4 = 8
  \label{eq:dim_total}
\end{equation}
real degrees of freedom.

\subsection{Gauge-fixing reduction}

On the winding saddle \eqref{eq:winding_saddle}, the gauge group acting on
$\delta\Theta$ is the stabiliser of the background, which contains at minimum
$U(1)_{\mathrm{EM}}$ (1 real parameter).  The full gauge group in the
alpha-sector derivation is $\mathrm{SU}(2)_L \times U(1)_Y$ (4 real parameters).

After gauge-fixing the $\mathrm{SU}(2)_L\times U(1)_Y$ symmetry (4 gauge
parameters), the number of \emph{physical} real d.o.f.\ of $\delta\Theta$ is:
\begin{equation}
  N_{\mathrm{phys}} \;=\; 8 - 4 \;=\; 4.
  \label{eq:N_phys}
\end{equation}

\subsection{Winding structure and winding-relevant d.o.f.}

The winding saddle $\bar\Theta_n = \Theta_0 e^{in\psi}$ is associated with
the $\psi$-circle winding.  Only the three imaginary quaternion components
$(\Theta_1, \Theta_2, \Theta_3)$ wind around the compact imaginary-time
directions (see \texttt{step1\_mode\_decomposition.tex} §2.2); the real
component $\Theta_0$ does not wind around $S^1_\psi$ in the EM sense.

This gives a winding-sector counting argument:

\begin{itemize}
  \item Physical d.o.f.\ from $\mathrm{Im}\,\HH$ after gauge-fixing:
    $(\Theta_1, \Theta_2, \Theta_3)$ are 3 complex fields = 6 real d.o.f.
    From these, 3 real Goldstone phases are eaten by
    $\mathrm{SU}(2)_L\times U(1)_Y$ (Higgs mechanism at the saddle).
    Remaining: $6 - 3 = 3$ real d.o.f.

  \item Alternatively: $N_{\mathrm{phys}} = 4$ (from above), minus 1 (the
    neutral $\Theta_0$ component that does not wind) $= 3$ winding d.o.f.
\end{itemize}

Both estimates give 3 physical real winding degrees of freedom contributing
to the 2D CFT on $S^1_t\times S^1_\psi$.

\begin{proposition}[Mode-counting argument for $c = 3$]
\label{prop:c3}
A mode-counting argument from $\delta^2S[\Theta]/\delta\Theta^2$ at the
winding saddle gives:
\begin{equation}
  c \;=\; N_{\mathrm{winding\,phys.}} \;=\; 3.
  \label{eq:c_3}
\end{equation}
\end{proposition}

%──────────────────────────────────────────────────────────────────────────────
\section{Step 4 — Gap Assessment}
\label{sec:gap}
%──────────────────────────────────────────────────────────────────────────────

\subsection{What is proved}

\begin{itemize}
  \item The fluctuation operator $\hat{K}_n = -\partial^2 + m_n^2$ is
        derived from $S_{\mathrm{kin}}[\Theta]$ at the winding saddle. [L1]
  \item In the massless limit, the 2D theory on $S^1_t\times S^1_\psi$
        is a free scalar CFT with $c = N_{\mathrm{dof}}$ per real scalar. [STD]
  \item A mode-counting argument gives $N_{\mathrm{dof}} = 3$. [MC]
\end{itemize}

\subsection{What is not proved}

\begin{gapbox}
\textbf{Gap G-c3 (partially narrowed):} The following steps remain unproved
at [L1] level:
\begin{enumerate}
  \item The exact gauge-fixing count: which of the 4 SU(2)$_L\times$U(1)$_Y$
        generators are actually broken at the winding saddle, and how many
        physical d.o.f.\ survive as winding modes.  The count $N_\mathrm{phys} = 3$
        relies on a specific symmetry-breaking pattern at the saddle.
  \item The identification of the resulting 2D theory as a \emph{unitary} CFT
        with a well-defined Virasoro central charge.  The free massive scalar
        is not a CFT (it has a mass gap); the identification requires a
        decoupling limit or an RG-flow argument.
  \item The derivation of $c$ from the OPE
        $T(z)T(0) \sim c/(2z^4) + \ldots$ rather than from mode-counting.
\end{enumerate}
\end{gapbox}

\subsection{Strengthened conclusion}

The mode-counting argument provides stronger support for $c = 3$ than was
available before: it derives $N_{\mathrm{dof}} = 3$ directly from counting
physical winding degrees of freedom at the winding saddle of $S[\Theta]$,
using the gauge structure of $\BB = \CC\otimes\HH$.

However, the derivation is \emph{not} at proof level [L1].  The status is:
\begin{itemize}
  \item \textbf{Before this document}: G-c3 OPEN (mode-counting argument only,
        no derivation from $\delta^2 S$).
  \item \textbf{After this document}: G-c3 CONDITIONAL [MC] — the
        mode-counting argument is now explicitly tied to the fluctuation
        operator $\hat K_n$ derived from $S[\Theta]$, and gives $c = 3$.
        A full [L1] proof requires the additional steps in §\ref{sec:gap}.
\end{itemize}

\begin{lemma}[Gap G-c3 current status]
\label{lem:c3}
From the quadratic expansion of $S_{\mathrm{kin}}[\Theta]$ around the winding
saddle $\bar\Theta_n = \Theta_0 e^{in\psi}$, a mode-counting argument gives:
\begin{equation}
  c \;=\; 3.
\end{equation}
This result is \textbf{CONDITIONAL [MC]}: it follows from the winding-mode
degree-of-freedom count after gauge fixing, but not yet from a complete
functional-determinant computation or OPE calculation.
\end{lemma}

\begin{statusbox}
Výsledek: $c = 3$ (mode-counting argument from $\delta^2 S[\Theta]/\delta\Theta^2$
at winding saddle).\\
Klasifikace: CONDITIONAL [MC] — not [L1] because gauge-fixing count is
  not fully rigorous and CFT identification requires RG-flow or decoupling argument.\\
Dopad na alpha track: If $c = 3$ is accepted, the exponent $c/12 = 1/4$
  in the formula $B = N_{\mathrm{eff}}^{3/2}(2\eta(i))^{c/12}$ acquires
  a mode-counting justification.  Gap G-c3 is narrowed from OPEN to
  CONDITIONAL [MC].  A full [L1] proof requires further work.
\end{statusbox}

\ifdefined\INCLUDEMODE
\else
  \end{document}
\fi
